\section{Introduction}\label{sec:intro}

When a trade shock destroys a job in Britain, the state's response is
improvised. The firm may be rescued---the interventions around the two
integrated steelworks have committed roughly
\pounds\InstAdHocSteelTotal\,million of public money since
2024---but the worker faces the ordinary benefit system: a contributory
allowance of \pounds\TaaJsaWeekly\ a week for six months, then
means-tested Universal Credit gated by savings rules and claiming
behaviour. The United States has operated statutory Trade Adjustment
Assistance in some form since 1962 and has trialled wage insurance for
older displaced workers; the European Union runs a Globalisation
Adjustment Fund, which the United Kingdom left at Brexit without
replacement. The sharper irony is domestic: the UK itself ran a statutory
wage-insurance scheme for roughly three decades---the Iron and Steel
Employees Re-adaptation Benefits Scheme paid weekly supplements topping
redundant steelworkers' new earnings up to 90 per cent of their previous
pay, part-funded by the European Coal and Steel Community---before letting
it lapse with the ECSC Treaty in 2002 \citep{iserbs1988}. Twenty years
later, the state is improvising rescue packages at the same steelworks the
scheme once covered. What no one has done is design and cost a
\emph{modern} UK instrument of this kind: the existing costed proposals
are economy-wide earnings-related unemployment insurance
\citep{resolutionfoundation2023, ifs2025ui, ippr2011}---paid while out of
work, not targeted at displacement, with no re-employment top-up---and
the trade-targeted wage-insurance evidence is entirely American. That is
the gap this paper fills.

The absence matters now for three reasons. First, the 2025 US tariff
schedule and its partial mitigation have made concentrated, trade-driven
job loss a live policy scenario rather than a historical one: the
companion study to this paper \citep{ahmadi2026} finds that a
tariff-calibrated earnings loss delivered through displacement is cushioned
by the existing tax--benefit system at roughly
\TaaCompanionDisplacedCushion\ per cent, against
\TaaCompanionWageCushion\ per cent when the same loss arrives as broad
wage cuts---and that the displacement-side cushion depends heavily on
Universal Credit claiming behaviour. The workers most exposed to a trade
shock are, in other words, the least insured, and their insurance is the
most fragile. Second, the government is already paying for the absence:
the ad-hoc steel packages are worker-and-community support improvised
deal by deal, uncosted as insurance, unevaluated, and unavailable to
workers displaced outside the headline interventions. Third, the
machinery for a designed alternative already exists inside the welfare
reform programme: the government has proposed replacing New Style
Jobseeker's Allowance and New Style Employment and Support Allowance with
a time-limited earnings-linked Unemployment Insurance---a reform whose
parameters can be applied directly to displaced workers and costed as one
arm of a trade-adjustment policy.

This paper does three things. First, it specifies four candidate
worker-side instruments precisely enough to cost: (i) \emph{wage
insurance} on the US Reemployment Trade Adjustment Assistance model---a
payment covering half the gap between pre- and post-displacement earnings
for up to two years of re-employment; (ii) the government's proposed
\emph{Unemployment Insurance} benefit, applied at its stated rate and
candidate durations to displaced workers as an incremental extension of
contributory support; (iii) a flat \emph{retraining stipend}; and (iv) a
\emph{redundancy-pay disregard} in the Universal Credit capital test, so
that statutory redundancy payments do not extinguish the means-tested
safety net they mechanically interact with. Second, it produces
first-pass annual costings of each arm from public data alone---ONS
redundancy flows, ASHE earnings, and statutory benefit parameters---under
a narrow eligibility rule (manufacturing redundancies) and a broad one
(all redundancies), with take-up and re-employment scenarios bracketed
rather than assumed away. Third, it specifies the microsimulation stage
that turns gross costings into net household and Exchequer incidence: the
same Family Resources Survey and PolicyEngine UK architecture as the
companion study, in which each arm's payments interact with income tax,
National Insurance and the Universal Credit taper, so that a pound of
gross wage insurance does not reach households as a pound of disposable
income.

Three design principles run through the exercise. The costings are
\emph{envelopes, not forecasts}: every behavioural margin---take-up of a
benefit that does not yet exist, re-employment responses to wage
insurance, employer responses to a disregard---is bracketed by scenario
rather than estimated, and the US evidence that wage insurance shortens
non-employment \citep{hyman2024} enters as a labelled sensitivity, not a
baseline assumption. Eligibility is \emph{bounded, not identified}:
``trade-displaced'' is not observable in UK survey data and would not be
administrable without a certification mechanism, so every arm is costed
under a narrow trade-proxy rule and a broad any-redundancy rule, the
policy-relevant range within which any real scheme would sit. And every
reported number is \emph{generated, not typed}: the costing pipeline
downloads its inputs, pins them, and emits the manuscript's values as
machine-generated macros, following the conventions of the companion
study's replication package.

The headline arithmetic is instructive even before the microsimulation
stage. Annual manufacturing redundancies run at roughly
\TaaFlowManuf\ workers a year; costing the
wage-insurance arm at central parameters puts its gross cost at roughly
\pounds\TaaWageInsCostCentralNarrow\,million a year under narrow eligibility.
The four arms together cost
\pounds\TaaTotalCostCentralNarrow\,million a year under narrow
eligibility---that is, a designed, rules-based, ex-ante insurance system
for displaced manufacturing workers would cost the Exchequer the same
order of magnitude as the improvised support already committed around two
steelworks (\TaaTotalCentralPctSteelNarrow\ per cent of it). The comparison is not decisive---the
steel packages buy industrial capacity as well as worker support---but it
reframes the question: the UK is already paying TAA-scale money for
TAA-shaped problems, without the design, the coverage, or the
evaluation that a statutory instrument would carry.

The rest of the paper proceeds as follows. Section~\ref{sec:literature}
locates the exercise in the US evaluation literature, the wage-insurance
design canon, and the UK evidence on displacement costs.
Section~\ref{sec:institutional} documents the institutional gap: the
ad-hoc steel interventions, the welfare-reform baseline, and the
trade-strategy asymmetry. Section~\ref{sec:design} specifies the four
arms. Section~\ref{sec:method} sets out the public-data costing method
and its deliberate limits. Section~\ref{sec:results} reports the
costings. Section~\ref{sec:roadmap} specifies the microsimulation stage.
Section~\ref{sec:discussion} concludes.
